% !TeX root = ../教材题库_全部条目_转写工作稿.tex
% Source: XJ-B1-C05-HW-P187-Y0986
% 第5章 / 5.3 三角函数的图象与性质 / 教材 p.187 / PDF 第194页
% OCR未校对：公式与文字均可能有误；未完成数学TeX转写；禁止直接发给学生
% 原题及相邻原解截图（核对用）：../题目原图/XJ-B1-C05-HW-P187-Y0986-01.webp
% 本文件只在转写完成后作为题干使用；OCR 条目尚不是可靠题干。
\begin{bankitem}{习题5.3 第12题}{第5章 / 5.3 三角函数的图象与性质 / 教材 p.187 / PDF 第194页}{XJ-B1-C05-HW-P187-Y0986}{OCR未校对：公式与文字均可能有误；未完成数学TeX转写；禁止直接发给学生}
\begin{ocrdraft}
[教材p.187 / PDF第194页]
\par
12．
\par
正弦函数是奇函数，它的图象关于原点对称，即原点是它的对称中心
\par
（1)除原点外，它还有其他的对称中心吗？若有，对称中心的坐标是什么？
\par
（2）正弦曲线是轴对称图形吗？如果是，对称轴的方程是什么？
\par
函数3一3，（2，十晋）的图象是中。对称图形吗？是轴对称图形吗？若是，
\par
（3）
\par
对称中心的坐标和对称轴的方程分别是什么？
\end{ocrdraft}
\end{bankitem}
